\documentclass[11pt]{article}
\usepackage{mathblatt}
\makeatletter
\setlength{\mbpfbe}{0mm}% keine Punkte (Bauregel 6.4)
\makeatother
% Rechenraum als Karo (Bauregel 6.9): n Rechenschritte = n mal 10 mm
\newcommand{\karo}[1]{\par\smallskip\noindent\begin{tikzpicture}
  \clip (0,0) rectangle ({\linewidth-1mm},{#1*10mm});
  \draw[black!16,line width=0.3pt,step=5mm] (0,0) grid ({\linewidth},{#1*10mm});
  \end{tikzpicture}\par}
\newcommand{\linie}[1]{\,{\color{black!45}\rule[-1pt]{#1}{0.4pt}}\,}
\newcommand{\pkt}{\textperiodcentered{}}
\begin{document}

\pfheftstil
\pflaufkopf[P10 \pkt{} mehr]{Seite oder Winkel mit sin, cos, tan}{B8X}\pfseitevon
\pfheftkopf{Seite oder Winkel mit sin, cos, tan}{Portion 2 von 2}
\pfstufenzeile{Seite berechnen \pkt{} Seitenverhältnis benennen \pkt{} Winkel berechnen}
% 2021-OS-K3a-m1 (Schwere I, rang 8)
\begin{pfaufg}{nach~P10\\’21~\pkt{}~3a}{1.}{}
Gegeben ist das Viereck ABCD mit der Diagonale BD = 1400 m. Bei A ist ein rechter Winkel. Die Seite DC ist 1950 m lang. Bei D ist der Winkel zwischen DA und DB 53° groß und der Winkel zwischen DB und DC 62°. Bei B ist der Innenwinkel $\angle$ABC = 112° groß. Zeige mit einer Rechnung, dass die Strecke AB etwa 1118 m lang ist.
\par\smallskip\noindent{\small\itshape Skizze: Viereck: A rechts unten, B senkrecht über A, D links von A, C weit links oben; Diagonale BD = 1400 m; rechter Winkel bei A; $\angle$ADB = 53°; nicht maßstabsgerecht}\par
\karo{2}
\fusshilfe{\mbox{1:~1118 m}}
\end{pfaufg}
% 2026-FOR-K4c-m1 (Schwere II, rang 6)
\begin{pfaufg}{nach~P10\\’26~\pkt{}~4c}{2.}{}
Im Viereck ABCD sind gegenüberliegende Seiten parallel. Die Seite BC ist 38 cm lang. Die Seite AB wird über B hinaus bis F verlängert; BF ist 16 cm lang. Die Höhe h = CF steht senkrecht auf AF. Der Winkel bei B zwischen BF und BC heißt $\varepsilon$. Die Diagonale AC bildet mit AB einen Winkel von 24°. Bestimme die Länge von AC. Gegeben: h $\approx$ 34,47 cm.
\par\smallskip\noindent{\small\itshape Skizze: Viereck (Parallelogramm): A links unten, B rechts unten, C rechts oben, D links oben, nach rechts geneigt; BC = 35 cm; AB bis F verlängert; h gestrichelt von C nach F, rechter Winkel bei F; Diagonale AC, Winkel 24° bei A zwischen AB und AC; nicht maßstabsgerecht}\par
\karo{3}
\fusshilfe{\mbox{2:~84,7 cm}}
\end{pfaufg}
% 2020-OS-B1c-m1 (Schwere I, rang 10)
\begin{pfaufg}{nach~P10\\’20~\pkt{}~1c}{3.}{}
Schreibe $\tan\alpha$ als Verhältnis zweier Seiten des abgebildeten Dreiecks.
\par\smallskip\noindent{\small\itshape Skizze: rechtwinkliges Dreieck: rechter Winkel an der oberen Ecke zwischen s (links) und r (rechts); Hypotenuse t unten; $\alpha$ unten links zwischen s und t}\par
\hfill tan $\alpha$ $=$\linie{30mm}
\fusshilfe{\mbox{3:~r/s}}
\end{pfaufg}
% 2019-OS-K3b-m1 (Schwere II, rang 16)
\begin{pfaufg}{nach~P10\\’19~\pkt{}~3b}{4.}{}
Beim Fußballtraining spielen sich drei Kinder den Ball zu: Amir steht bei A, Ben bei B, Chris bei C. Zwischen Ben und Amir sind es 5 m, zwischen Amir und Chris 12 m. Bei B ist ein rechter Winkel; der Winkel bei A heißt $\alpha$. Zeige mit einer Rechnung, dass $\alpha$ etwa 65° groß ist.
\par\smallskip\noindent{\small\itshape Skizze: rechtwinkliges Dreieck: rechter Winkel bei B (unten links), A senkrecht über B, C rechts von B; AB = 5 m, AC = 12 m; $\alpha$ bei A; nicht maßstabsgerecht}\par
\karo{3}
\fusshilfe{\mbox{4:~65°}}
\end{pfaufg}
% 2019-OS-B1h-m1 (Schwere I, rang 12)
\begin{pfaufg}{nach~P10\\’19~\pkt{}~1h}{5.}{}
Schreibe $\sin\alpha$ als Verhältnis zweier Seiten des abgebildeten Dreiecks.
\par\smallskip\noindent{\small\itshape Skizze: rechtwinkliges Dreieck: rechter Winkel an der oberen Ecke; Hypotenuse t waagerecht unten; Seite s von links unten zur Spitze, Seite r von der Spitze nach rechts unten; Winkel $\alpha$ unten links zwischen s und t}\par
\hfill sin $\alpha$ $=$\linie{30mm}
\fusshilfe{\mbox{5:~r/t}}
\end{pfaufg}
% 2017-OS-B1j-m1 (Schwere I, rang 14)
\begin{pfaufg}{nach~P10\\’17~\pkt{}~1j}{6.}{}
Schreibe $\sin\beta$ als Verhältnis zweier Seiten des abgebildeten Dreiecks.
\par\smallskip\noindent{\small\itshape Skizze: rechtwinkliges Dreieck: rechter Winkel unten links; senkrechte Kathete r, waagerechte Kathete s, Hypotenuse t; Winkel $\beta$ unten rechts}\par
\hfill sin $\beta$ $=$\linie{30mm}
\fusshilfe{\mbox{6:~r/t}}
\end{pfaufg}
% 2017-OS-K4b-m1 (Schwere II, rang 18)
\begin{pfaufg}{nach~P10\\’17~\pkt{}~4b}{7.}{}
Ein Dachboden hat als Querschnitt das Dreieck ABC. Die Dachschräge AC ist 4,90 m lang. Der Winkel bei A beträgt 64,4°, der Winkel bei B 36,1°; der Winkel bei C heißt $\gamma$. Auf dem Boden AB steht im Punkt P eine senkrechte Wand, die genau bis zur Dachschräge BC reicht. Sie ist 1,35 m hoch. Berechne den Abstand von P bis B.
\par\smallskip\noindent{\small\itshape Skizze: Dreieck: A links unten, B rechts unten, C oben links von der Mitte; 36,2° bei B; senkrechte Wand von P auf AB (nahe B) bis zur Schräge BC, 1,35 m hoch; nicht maßstabsgerecht}\par
\karo{2}
\fusshilfe{\mbox{7:~1,85 m}}
\end{pfaufg}
\par\vfill\noindent{\scriptsize\color{mbgrau}Originale: P10 2021 \pkt{} 3a, P10 2026 \pkt{} 4c, P10 2020 \pkt{} 1c, P10 2019 \pkt{} 3b, P10 2019 \pkt{} 1h, P10 2017 \pkt{} 1j, P10 2017 \pkt{} 4b\par}
\end{document}